\scanpage{271}

equivalence \inserted{\unsure{of} 2-categories,} i.e. that it is \unsure{3} faithful.

Of course we are reduced all the \unsure{same} to the case of groupoids \inserted{\unsure{``connected''}} \(\Pi_f\) defined by a group \inserted{(\unsure{finite} if need be)}, and we are also reduced by \unsure{essentially} trivial arguments to looking at the case of diagrams \(B_{\unsure{D^*}}\to B_U\) with \(B_U\) connected.

On the geometric side, the situation \unsure{would be} given by a \(U=X\setminus S\) of type \((g,\nu)\) (\((g,\nu)\ne(0,0)\)), and an operation of \(\Gamma\) on it --- we may suppose \(U=U_{g,\nu}\), so we give \(\Gamma\to A_{g,\nu}\). For the question of \unsure{\(i\)-faithfulness} with \(i\leq 2\), we may suppose that it is the \unsure{same} group \(\Gamma\) that operates\unsure{;} \deleted{\illegible{4}}

\unsure{1)} \uline{0-faithfulness}. \deleted{\illegible{1}} Let \(U,U'\) be given with an operation of \(\Gamma\) on them, \(f,g:U\stackrel{\sim}{\longrightarrow}U'\) commuting with \(\Gamma\), and \(\alpha,\beta\) two homotopies \inserted{\unsure{from \(f\) to \(g\)}} i.e. two paths from \(f\) to \(g\) in \(\uline{\mathrm{Isom}}_\Gamma(U,U')\), we suppose that in \unsure{the topossic description|the topossic descriptions}
\diagram{\begin{tikzcd}[ampersand replacement=\&,column sep=5.5em,row sep=3.5em]
B_{D^*,\Gamma}
 \arrow[r]
 \arrow[d,bend right=16,"f_{D^*,\Gamma}"']
 \arrow[d,bend left=16,"g_{D^*,\Gamma}"]
\& B_{U,\Gamma}
 \arrow[r,no head]
 \arrow[d,bend right=16,"f_{U,\Gamma}"']
 \arrow[d,bend left=16,"g_{U,\Gamma}"]
\& B_\Gamma \arrow[d,equal]
\\
\unsure{B_{D'^*,\Gamma}} \arrow[r,no head]
\& B_{U',\Gamma} \arrow[r,no head]
\& B_\Gamma
\end{tikzcd}}
we have \(\alpha_*^{D^*}=\beta_*^{D^*}:f_{D^*}\to g_{D^*}\) and \(\alpha_*^U=\beta_*^U:f_U\to g_U\). To prove that \(\alpha\) and \(\beta\) are homotopic.

\margin{270}

\scanpage{272}
\mspage{137}

\margin{i.e. \(f\in A_{g,\nu}\) commuting with the image of \(\Gamma\) under \(\Gamma\to A_{g,\nu}\).}
We may suppose \(g=\mathrm{id}\), \(U=U_{g,\nu}\), so \(f\in A_{g,\nu}^{\Gamma}\), and we have two paths \(\alpha,\beta\) from \(1\) to \deleted{\illegible{1}} \(f\) in the space \(A_{g,\nu}^{\Gamma}\). We suppose that the two corresponding isomorphisms between the identity \inserted{of \unsure{\(\Pi_{D^*}\)}} and \unsure{\(f_{D^*,\Gamma}\)} : \(\Pi_{\deleted{\unsure{B}}\unsure{D^*_{S_\nu},\Gamma}}\to\Pi_{\unsure{D^*_{S_\nu},\Gamma}}\) on the one hand, between the identity of \(\Pi_{U,\Gamma}\) and \(f_{U,\Gamma}\) on the other hand, are the same, we want to prove that \(\alpha\) and \(\beta\) are homotopic. Of course, the hypothesis \deleted{\illegible{1}} on \(\alpha,\beta\) concerning the action of \(\Gamma\) is satisfied a fortiori upon restricting to a smaller group, for example \deleted{\unsure{\(\Gamma\)}} \(\Gamma'=1\), and in fact we see that it is equivalent for \(\Gamma\), and for its subgroup \(1\). \deleted{\illegible{1}} the already known result (\unsure{Ladegallerie|Ladegaillerie}) for \(\Gamma=1\), shows that this means that if two paths \inserted{in \(A^{\Gamma}_{\deleted{\unsure{g,\nu}}}\)} from \(1\) to \(f\) are homotopic in \(A\), they are so in \(A^{\Gamma}\) or again that
\margin{\unsure{let us put} \(A=A_{g,\nu}\)}
\[
\fbox{\(\pi_1(A^{\Gamma})\to\pi_1(A)\ \mbox{is injective.}\)}
\]
In the case \(\pi_1(A)=0\) (case \((g,\nu)\) anabelian) this thus comes down to proving that \(\pi_1(A^{\Gamma})=0\).

\margin{271}

\scanpage{273}

2) \uline{1-faithfulness}. This means \inserted{(in addition to 0-faithfulness)} that \unsure{\(\forall\)} isomorphism between \deleted{\unsure{\(f\)}} \(B_f\) and \(B_g\) comes from a \deleted{\unsure{path}} path from \(f\) to \(g\). With \unsure{the preceding reductions|the preceding reduction}, we may suppose \(g=\mathrm{id}\), so we have \(f\in A^{\Gamma}\), whence \unsure{\(f^*:(B_{D^*}\to B_U)\to(B_{D^*}\to B_U)\)} (\unsure{commuting} \deleted{\unsure{with}} \(\Gamma\) \inserted{\unsure{respecting \(\Gamma\)}}), and we have an isomorphism with the identity --- unless I am mistaken this condition

\margin{272}

\scanpage{274}
\mspage{138}

Let \(A\) be a topological group, whence two invariants
\[
\pi_0=\pi_0(A),\qquad \pi_1=\pi_1(A)
\]
the first a group \inserted{not necessarily commutative}, \deleted{\illegible{1}} the second a commutative group on which \(\pi_0\) operates. We define a Gr-category \inserted{\(\uline{A}\)} with invariants \(\pi_0,\pi_1\) and corresponding to this operation of \(\pi_0\) on \(\pi_1\), by taking as underlying category the \deleted{\illegible{3}} fundamental \deleted{\unsure{Hom}} groupoid (naive) of \(A\), and as composition functor the one deduced from \(A\times A\to A\) (associativity in \(\uline{A}=\Pi A\) is strict \ldots)

This being so, every torsor under \(A\) defines a 1-torsor under the Gr-category \(\uline{A}\). More generally, for \(\forall\) topological space \(S\) (or \(\forall\) topos which is locally a topological space) every torsor on \(S\) with group \(A_S\) defines an \inserted{\(\uline{A}_S\)-}stack on \(S\) which is a \deleted{\unsure{\(\uline{A}_S\)}} stack of \(A_S\)-torsors. If \(\pi_i(A)=0\) for \(i\geq2\), then we find

\margin{273}

\scanpage{275}

thus \inserted{for \unsure{\(\forall n\in N,\ n\geq2\)}} an \unsure{\(n\)}-\uline{equivalence} between the \inserted{\(n\)-}category of torsors on \(S\) with group \(A_S\), and the \unsure{\(n\)}-category deduced from the 2-category of \inserted{\unsure{gerbes}} \(\uline{A}_S\)-torsors by extending \unsure{it} in a discrete way \ldots

But let us consider \deleted{\illegible{1}} a group \(\Gamma\), and consider the classification of \(A_{B_\Gamma}\)-torsors on the classifying topos --- i.e. that of \(A\) torsors, with an operation of \(\Gamma\) on them (commuting with the action of \(A\)). These objects form a 2-category, whose connected components correspond to the conjugacy classes of homomorphisms from \(\Gamma\) to \(A\). If we have \deleted{two} \inserted{one} homomorphism \(\Gamma\to A\), it defines \deleted{an} an action on the trivial torsor. \(1_A\) If \(u,v\) are two homomorphisms, the isomorphisms from \((1_A,u)\) onto \((1_A,v)\) correspond to the elements of the set \unsure{``}
\[
\unsure{\mathrm{Transp}_A}(u,v)=\{g\in A\mid v=\mathrm{int}(g).u\}
\]
But we shall make a 0-isotopic category,

\margin{274}

\scanpage{276}
\mspage{139}

by replacing \(\mathrm{Transp}(u,v)\) by its \(\pi_0\), and a 1-isotopic category by replacing \inserted{it} by its fundamental groupoid \(\Pi_1\mathrm{Transp}(u,v)\). Thus an isomorphism from \(u\) to \(v\) is still a point of \(\mathrm{Transp}(u,v)\), but if we have two such isomorphisms \inserted{\(f,g\)}, the isomorphisms \(f\simeq g\) are the classes of \unsure{paths} in \(\mathrm{Transp}(u,v)\) from \(f\) to \(g\). Thus, the category \(\uline{\mathrm{Aut}}(u)\) is equivalent to \(\Pi_1\underbrace{\unsure{\mathrm{Trans}}(u,u)}_{A^u}\) \deleted{is}
\[
\uline{\mathrm{Aut}}(u)\simeq\Pi_1\unsure{\mathrm{Trans}}(u)
\]
and the \deleted{\unsure{the}} category \(\mathrm{Isom}(u,v)\) is either empty (if \(u,v\) are not conjugate), or a torsor under \(\uline{\mathrm{Aut}}(u)\).

But for every torsor \inserted{\(P\)} over \(A\) on which \(\Gamma\) operates, \deleted{\illegible{1}}\inserted{the} \deleted{\unsure{associated a}} 1-torsor \deleted{under} \(\Pi_1P\) under \(\Pi_1A\) \inserted{\(=\uline{A}\)} is equipped with operations of \(\Gamma\), whence a

\margin{275}

\scanpage{277}

\unsure{\((\Gamma,-\uline{A})\)} torsor. We thus find a functor of 2-categories
\diagram{\begin{tikzcd}[ampersand replacement=\&,column sep=large]
{\begin{array}{l}
\mbox{2 category of }\deleted{\mbox{torsors}}\\
\Gamma-A\mbox{-torsors}\\
\mbox{[or again, of homomorphisms }\Gamma\to A\mbox{]}
\end{array}}
\arrow[r] \&
{\begin{array}{l}
\mbox{2 category of 1-torsors under }A,\\
\mbox{with operations of }\Gamma\mbox{ on them.}
\end{array}}
\end{tikzcd}}

In short, we have just repeated on the classifying topos \(B_\Gamma\) the construction made \inserted{above} for a topological space \(S\). We would also like to \inserted{express} conditions for this \deleted{homomorphism} \inserted{2-functor} to be a 2-equivalence. For this, we should first have an understanding of the classification of the classes \unsure{of equivalence|of equivalences} of objects of the second category --- which are, for their part, of a purely algebraic nature, in terms of the Gr-category \(\uline{A}\) and of \(\Gamma\). \unsure{Even if} we restrict ourselves to

\margin{276}

\scanpage{278}
\mspage{140}

trivial torsors over \(\uline{A}\) (which is legitimate) we would have to make explicit what it means that \(\Gamma\) operates on the trivial torsor. We observe that this means that we have a homomorphism \inserted{of Gr-categories} from \inserted{the discrete Gr-category defined by} \(\Gamma\) to \(\uline{A}\). Thus the 2-category envisaged is the one whose objects are the homomorphisms of Gr-categories
\diagram{\begin{tikzcd}[ampersand replacement=\&]
\Gamma \arrow[r,"u","{\deleted{\unsure{\underset{v}{\longrightarrow}}}}"'] \& \uline{A}
\end{tikzcd}}
and for two such morphisms \(u,v\), we must define (not only a set \(\mathrm{Hom}(u,v)\) but) a category \(\uline{\mathrm{Hom}}(u,v)\), as the subcategory \(\uline{\mathrm{Transp}}(u,v)\) of \(\uline{A}\), to be defined ad-hoc.

Let us consider for simplicity the case where \(\pi_1=0\), then \(\uline{A}\) reduces to \(\pi_0\), and the category of homomorphisms \(\Gamma\to\uline{A}\) to that of homomorphisms \(\Gamma\to\pi_0\). The

\margin{277}

\scanpage{279}

fact that the functor of 2-categories above \deleted{\unsure{that it}} is an equivalence of \unsure{2|1}-categories then translates step by step as follows :

1°) \uline{0 faithfulness} means that for \(\forall\) homomorphism \(u:\Gamma\to A\), we have \(\pi_1(\unsure{A^u})=0\)

\inserted{\uline{N.B.} \(A^u\) and \unsure{\(A^{0u}\)} have the same \unsure{identity component|identity components}, so the condition is written : \fbox{\(\pi_1(\unsure{(A^0)^u})=0\)}.}

2°) \uline{1-faithfulness} means that moreover, if \(u,v\) are \deleted{two} homomorphisms \(\Gamma\rightrightarrows A\), and \deleted{\unsure{\(\alpha,\beta\in\mathrm{Transp}\)} \illegible{1}} \(\alpha,\beta\in\mathrm{Transp}_A(u,v)\) \deleted{\illegible{1} by \illegible{2}} \inserted{have the same image in \(\mathrm{Transp}_{\pi_0}(\uline{u},\uline{v})\)} \deleted{then \illegible{1} \(\alpha,\beta\) \illegible{3} \(A\to\pi_0\) are conjugate, so \unsure{\((u,v)\)} \illegible{1}} \deleted{\illegible{1}} then \(\alpha,\beta\) are in the same \deleted{\illegible{1}, even \unsure{identical}), then \(u\) and \(v\) are} connected component of \(\mathrm{Transp}_A(u,v)\). \deleted{\unsure{This}} \deleted{\unsure{conjugate to one another}} \deleted{\unsure{To}} To see this, we may suppose \(u=v\), and \margin{\unsure{3°)}} we are reduced to expressing that the map \(\pi_0(A^u)\to\pi_0\) is injective i.e. that \deleted{\unsure{\(A^u\)}} its kernel is reduced to \(1\), i.e. that \inserted{the open subgroup} \(\unsure{(A^u)\cap A^0=(A^0)^u}\) of \(A^u\) is connected, i.e.
\[
\fbox{\(\pi_0(\unsure{(A^0)^u})=0\)}
\]

\margin{278}

\scanpage{280}
\mspage{141}

3°) \uline{2-faithfulness} means that \unsure{in} addition to the preceding conditions, which ensure that for \(u,v\) fixed, the functor \(\uline{\mathrm{Hom}}(u,v)\to\uline{\mathrm{Hom}}(\uline{u},\uline{v})\) is fully faithful, that the latter is essentially surjective, i.e. that \unsure{if} \(u,v\) are such that \(\uline{u},\uline{v}:\Gamma\to\pi_0\) are conjugate, \deleted{then} (\unsure{equal if we wish}) \(u,v\) are already conjugate \inserted{\unsure{(by an element\(^{\mathrm t}\) of \(\uline{\pi_0}\))}}. In other words : for \(\forall\) homomorphism \(\uline{u}:\Gamma\to\pi_0=\unsure{A/A^0}\), if there exists a lifting of \(u\) to \(\Gamma\to A\), it is unique up to conjugation.

\margin{\unsure{We must say it} \illegible{1} \unsure{in a} \unsure{stronger way} : \(\mathrm{Transp}(u,v)\to\mathrm{Transp}(\uline{u},\uline{v})\) surjective --- this is equivalent to saying that if \unsure{\(\uline{u}=\uline{v}\)}, then \(u\) and \(v\) are conjugate by an \uline{element} \uline{of \unsure{\(A^0\)}}.}

4°) \uline{3-faithfulness} means that in addition, \unsure{every} homomorphism \(\uline{u}:\Gamma\to\pi_0\) lifts to \(u:\Gamma\to A\).

In summary, \inserted{if \(\pi_1(\unsure{A^0})=0\),} 3-faithfulness means that for \(\forall\) homomorphism \(\uline{u}:\Gamma\to\pi_0=\unsure{A/A^0}\), there exists \deleted{a} lifting \(u:\Gamma\to A\), unique up to conjugation, and that then \(\unsure{(A^0)^u}\) is connected and simply connected,

\margin{279}
